\chapter{TensorFlow, convolution, and observable training}
\label{ch:tensorflow-convolution}

\begin{learningobjectives}
After this chapter, you should be able to:
\begin{itemize}
  \item distinguish tensors, variables, automatic differentiation, Keras, and
  TensorBoard as separate parts of a learning system;
  \item translate a scientific modeling contract into the Keras model
  lifecycle;
  \item explain mathematically why convolution respects image structure;
  \item calculate convolution output shapes and receptive fields;
  \item distinguish the full-data MNIST experiment from its bounded,
  offline CI smoke test and design a defensible cats-versus-dogs study; and
  \item select optimization, generalization, data, and system telemetry that
  makes training diagnosable.
\end{itemize}
\end{learningobjectives}

\begin{chapterconnection}
Chapter~\ref{ch:neural-networks} constructed a neural network from literal
neurons and derived reverse-mode differentiation. The preceding MNIST lesson
then implemented the same mechanism directly. We now cross the framework
boundary. TensorFlow automates tensor execution and derivatives, while Keras
provides a higher-level model lifecycle. Neither removes the need to specify
data provenance, output semantics, held-out evidence, or operational
diagnostics.
\end{chapterconnection}

\section{A framework is a stack of contracts}

TensorFlow was released as open-source software in November 2015. Its original
white paper described both an interface for expressing machine-learning
algorithms and an implementation capable of executing computations across
heterogeneous systems.\footnote{M. Abadi et al., \emph{TensorFlow: Large-Scale
Machine Learning on Heterogeneous Distributed Systems}, 2015,
\url{https://www.tensorflow.org/static/extras/tensorflow-whitepaper2015.pdf}.}
TensorFlow~2 made eager execution the normal interactive experience and
standardized on Keras as its high-level interface. Modern Keras is a
multi-backend deep-learning API; this chapter uses its TensorFlow backend.

The names are related but not interchangeable:

\begin{description}
  \item[Tensor.] A typed multidimensional value with a shape and an execution
  device. A tensor is immutable.
  \item[Variable.] Mutable state, ordinarily a model parameter, tracked for
  differentiation and checkpointing.
  \item[Operation.] A transformation such as matrix multiplication,
  convolution, reduction, or activation.
  \item[Automatic differentiation.] A system that records differentiable
  operations and applies their local derivative rules in reverse.
  \item[Keras layer.] A reusable parameterized transformation with explicit
  input and output shapes.
  \item[Keras model.] A connected collection of layers with training,
  evaluation, inference, and serialization interfaces.
  \item[Callback.] A hook invoked at defined lifecycle events, for example to
  write telemetry, save a checkpoint, or stop training.
  \item[TensorBoard.] A separate application that reads event files and
  displays scalar series, histograms, images, graphs, and performance traces.
\end{description}

Figure~\ref{fig:tensorflow-stack} places these responsibilities in one system.

\begin{figure}[htbp]
\centering
\begin{tikzpicture}[
  node distance=4mm,
  box/.style={draw=RiceBlue, rounded corners=2mm, minimum width=10.4cm,
    minimum height=9mm, align=center, font=\small},
  arrow/.style={-{Latex[length=2.2mm]}, thick, RiceBlue}
]
\node[box, fill=RiceBlue!10] (question) {Scientific question, target,
  sampling process, split contract, decision cost};
\node[box, fill=RiceGray!16, below=of question] (keras) {Keras API: layers,
  models, losses, optimizers, metrics, callbacks};
\node[box, fill=RiceBlue!7, below=of keras] (tf) {TensorFlow runtime:
  tensors, variables, automatic differentiation, compiled device kernels};
\node[box, fill=RiceGray!16, below=of tf] (hardware) {CPU, GPU, accelerator,
  storage, and distributed execution};
\node[box, fill=RiceGold!18, below=of hardware]
  (observe) {Observability: Python logs, TensorBoard events, checkpoints, and
  artifact identity};
\draw[arrow] (question) -- (keras);
\draw[arrow] (keras) -- (tf);
\draw[arrow] (tf) -- (hardware);
\draw[arrow] (hardware) -- (observe);
\draw[arrow, dashed] (observe.east) to[bend left=38] (question.east);
\end{tikzpicture}
\caption{The framework occupies the computational middle of a larger
scientific system. It can calculate derivatives and execute kernels; it cannot
make the target valid or the evidence relevant. Observability connects the
running computation to human diagnosis.}
\label{fig:tensorflow-stack}
\end{figure}

\begin{misconceptionbox}
``Keras trained the model'' is not a complete experimental description. A
reproducible claim also needs the data version and sampling unit, split policy,
preprocessing fit boundary, architecture, initializers, objective, optimizer,
randomness policy, stopping rule, checkpoint-selection rule, metrics, software
versions, and final artifact identity.
\end{misconceptionbox}

\section{Tensors and automatic differentiation}

A batch of grayscale images commonly has shape
\[
  (n,h,w,1),
\]
where $n$ is batch size, $h$ and $w$ are spatial dimensions, and the last
coordinate is the channel. A color image normally has three input channels.
The channel dimension is semantic; silently transposing channel and spatial
axes changes the represented computation.

Let $L(\theta)$ be a scalar loss. Inside a \texttt{GradientTape}, TensorFlow
records relevant primitive operations. A call to \texttt{gradient} traverses
that recorded graph backward to compute $\nabla_\theta L$. This is the same
chain rule derived in Chapter~\ref{ch:neural-networks}. The framework maintains
the graph and supplies derivative implementations for supported operations.

\begin{definitionbox}{automatic differentiation is not symbolic or numerical differentiation}
Automatic differentiation evaluates exact chain-rule compositions of the
program's primitive derivative rules, up to floating-point arithmetic. It
does not manipulate a symbolic expression into a closed form, and it does not
approximate every coordinate by finite differences. Finite differences remain
valuable as an independent check on a small deterministic calculation.
\end{definitionbox}

The tape ordinarily watches trainable variables automatically. Constants must
be watched explicitly when derivatives with respect to them are required.
Integer-valued computations are not differentiable in this sense, stateful
updates can stop a gradient path, and persistent tapes retain more intermediate
state.\footnote{TensorFlow, ``Introduction to gradients and automatic
differentiation,'' \url{https://www.tensorflow.org/guide/autodiff}.}

\conceptcheck{A tape returns \texttt{None} for one parameter. Give three
distinct hypotheses: one involving connectivity, one involving dtype, and one
involving state. What evidence would distinguish them?}

\section{The Keras lifecycle is an executable argument}

A routine experiment passes through six stages:
\[
 \texttt{build}\longrightarrow\texttt{compile}\longrightarrow
 \texttt{fit}\longrightarrow\texttt{evaluate}\longrightarrow
 \texttt{predict}\longrightarrow\texttt{save}.
\]
Each verb answers a different question.

\begin{enumerate}
  \item \textbf{Build} establishes transformations, shapes, parameters, and
  output semantics.
  \item \textbf{Compile} binds an optimizer, loss, and reported metrics. These
  are not interchangeable: the loss supplies gradients, whereas a metric may
  be nondifferentiable.
  \item \textbf{Fit} updates parameters using training batches. Validation
  data may guide stopping or checkpoint selection but receives no gradient
  update.
  \item \textbf{Evaluate} computes declared metrics without parameter
  updates. The final test partition is used once after choices are frozen.
  \item \textbf{Predict} runs inference and must reproduce the training-time
  feature and preprocessing contract.
  \item \textbf{Save} preserves the chosen artifact. A fresh-process reload
  test verifies that serialization is operational, not merely promised.
\end{enumerate}

Keras supports a simple \texttt{Sequential} stack, a Functional API for graphs
with branches or shared layers, and model subclassing for custom behavior. Use
the least flexible interface that expresses the intended graph clearly.
Progressive disclosure is a virtue: a three-layer sequential model should not
be hidden inside a custom training architecture without a reason.

\section{MNIST: concise code, explicit semantics}

MNIST observations are $28\times28$ grayscale images of handwritten digits.
A dense baseline reshapes each image into
\[
  x\in\mathbb{R}^{784}.
\]
Flattening changes the view, not the values. It also erases the explicit
two-dimensional neighborhood relation. Dividing integer intensities by $255$
maps them to $[0,1]$ and makes the numerical range explicit.

For ten mutually exclusive classes, the model emits logits
$z\in\mathbb{R}^{10}$. With class index $y$, sparse categorical
cross-entropy is
\[
  \ell(z,y)=-z_y+\log\left(\sum_{k=0}^{9}\exp z_k\right).
\]
With a one-hot vector $e_y$, the same quantity can be written
\[
  -\sum_{k=0}^{9}(e_y)_k\log p_k,
  \qquad
  p_k=\frac{\exp z_k}{\sum_r\exp z_r}.
\]
Sparse and one-hot targets encode the same class but require different loss
interfaces. If a model emits raw logits, the loss must be told that its inputs
are logits. Applying softmax in the model and then configuring a logits loss
applies the wrong contract.

\subsection{Acquire data without hiding provenance}

The executable lesson has two routes with different evidential purposes.
During normal student work,
\texttt{tf.keras.datasets.mnist.load\_data()} obtains the full official
training and test arrays. On the first call, Keras downloads its MNIST archive
to the user's Keras cache; later calls ordinarily reuse the cached file. A
cache improves performance and supports repeated local work, but it is not a
complete provenance record. A professional experiment records the source,
version or checksum, license, sampling unit, and meaning of the supplied split.

Continuous integration cannot safely depend on network availability or repeat
a full scientific experiment for every commit. It therefore sets a documented
environment flag and loads a small, versioned repository fixture made from
authentic MNIST observations. This route exercises loading, representation
checks, model construction, fitting, evaluation, and telemetry. Its metrics
are \emph{software evidence only}. They must not be presented as the empirical
result of the full-data lesson.

\begin{definitionbox}{experiment versus smoke test}
An experiment estimates or compares scientifically relevant quantities under
a declared design. A smoke test asks whether a bounded execution path still
runs and satisfies basic contracts. Using real observations makes a smoke test
representative of data shape and dtype; it does not give the sample the scope
or statistical meaning of the full experiment.
\end{definitionbox}

\subsection{Validate first, then transform}

The acquired image arrays have shape $(n,28,28)$ and unsigned eight-bit pixel
values in $[0,255]$. Before training, the notebook checks image rank and shape,
label alignment, class vocabulary, and dtype. It then converts pixels to
\texttt{float32} and divides by $255$, producing values in $[0,1]$. This scale
factor is a fixed representation convention; unlike a fitted mean or standard
deviation, it does not estimate state from either partition.

The flattening operation remains inside the model. Thus an input batch with
shape $(n,28,28)$ becomes $(n,784)$ at a named model layer, and the serialized
model retains that shape transformation. Data validation belongs at the
boundary, while deterministic transformations may belong in the model artifact
when they must accompany inference.

\begin{workedexample}{a full-data MNIST baseline with a bounded CI analogue}
On the normal route, reserve the official test partition before modeling. Fit a
model containing \texttt{Flatten}, a 64-unit ReLU layer, and a ten-logit
output. The parameter count is
\[
  (784\cdot64+64)+(64\cdot10+10)=50{,}890.
\]
Use a declared seed, Adam with learning rate $10^{-3}$, sparse categorical
cross-entropy, and a maximum of two epochs for the live demonstration. Keras
reserves the last ten percent of the supplied training array as validation
data. This is transparent for the lesson; a research system should persist an
explicit split, especially when observations contain repeated subjects,
groups, or time order.

After architecture and training choices are frozen, report one official-test
loss and accuracy together with a row-true, column-predicted confusion matrix.
During CI, run the same interfaces against the bounded authentic-MNIST fixture
only to verify that the pipeline remains executable.

The result is a baseline, not a deployment claim. The experiment must also
record data source, counts, preprocessing range, parameter count, runtime, and
artifact identity. A bounded epoch budget is preferable in lecture; a longer
study belongs in an explicitly configured experiment, not in an import or
notebook-open side effect. CI remains bounded in observations as well as
epochs.
\end{workedexample}

\subsection{The API lifecycle marks changes of state}

Building establishes the forward computation and shapes. Compiling binds the
optimizer, differentiable objective, and reported metrics but changes no
parameter. Fitting updates parameters using training batches and evaluates the
declared validation partition after each epoch. The returned
\texttt{History} is in-memory telemetry, not the trained artifact itself.

Evaluation runs a frozen model without updates. Prediction returns ten logits
per observation; an \texttt{argmax} produces the displayed class. If an
application requires probabilities, it must apply the declared link function
and test calibration rather than calling arbitrary scores probabilities. The
confusion matrix then localizes which true digits become which predicted
digits, exposing error structure hidden by aggregate accuracy.

\section{Convolution preserves a spatial hypothesis}

A fully connected unit has an independent weight for every incoming pixel. A
convolution instead assumes that useful local patterns may recur at different
locations. It uses the same small kernel across the grid. This combines two
inductive biases:

\begin{itemize}
  \item \emph{local connectivity}: one activation initially sees a small
  neighborhood; and
  \item \emph{parameter sharing}: the same detector is applied at every
  eligible location.
\end{itemize}

For one input channel and kernel $K\in\mathbb{R}^{k_h\times k_w}$, the
operation called convolution in deep-learning libraries is ordinarily
cross-correlation:
\[
  z_{i,j}=b+
  \sum_{u=0}^{k_h-1}\sum_{v=0}^{k_w-1}
  K_{u,v}x_{i+u,j+v}.
\]
The mathematical convolution would reverse the kernel. The naming convention
does not impair learning because the kernel entries are learned, but the
distinction matters when reproducing a hand calculation.

For input size $n$, symmetric padding $p$, kernel size $k$, and stride $s$, an
output side has length
\[
  n_{\mathrm{out}}
  =\left\lfloor\frac{n+2p-k}{s}\right\rfloor+1.
\]
``Valid'' padding means $p=0$. ``Same'' padding is chosen to preserve spatial
size for stride one. A stride larger than one skips positions and reduces
resolution.

\begin{figure}[htbp]
\centering
\begin{tikzpicture}[
  pixel/.style={draw=RiceGray!65, minimum size=5.8mm, inner sep=0pt},
  kernel/.style={draw=RiceBlue, very thick, fill=RiceBlue!10, opacity=.72},
  arrow/.style={-{Latex[length=2mm]}, thick, RiceBlue}
]
\foreach \r in {0,...,4}{
  \foreach \c in {0,...,4}{
    \node[pixel] at (0.6*\c,-0.6*\r) {};
  }
}
\draw[kernel] (-0.3,0.3) rectangle (1.5,-1.5);
\node[above=4mm] at (1.2,0.3) {input patch $3\times3$};
\node[draw=RiceBlue, fill=RiceGold!18, rounded corners, minimum width=2.1cm,
  minimum height=1.1cm, align=center] (sum) at (4.5,-.6)
  {multiply\\and sum};
\draw[arrow] (1.65,-.6) -- (sum.west);
\node[draw=RiceBlue, fill=RiceBlue!10, minimum size=6mm] (out) at (7,-.6)
  {$z_{i,j}$};
\draw[arrow] (sum.east) -- (out.west);
\end{tikzpicture}
\caption{One convolution location. Every output value is an inner product
between a local patch and the same learned kernel, plus a bias.}
\label{fig:convolution-mechanics}
\end{figure}

\subsection{Channels and filters}

For $c_{\mathrm{in}}$ input channels, a filter has shape
$k_h\times k_w\times c_{\mathrm{in}}$. It spans every input channel and
produces one output channel. With $c_{\mathrm{out}}$ filters, the kernel tensor
contains
\[
  k_hk_wc_{\mathrm{in}}c_{\mathrm{out}}
\]
weights, plus $c_{\mathrm{out}}$ biases. The output channels are learned feature
maps, not copies of RGB.

The notebook's inspectable CNN therefore follows
\[
 (28,28,1)\to(28,28,16)\to(14,14,16)
 \to(14,14,32)\to(7,7,32)\to32\to10.
\]
The two max-pooling layers have no trainable parameters, and global average
pooling replaces each $7\times7$ feature map by its mean. The parameter count is
\[
  16(3\cdot3\cdot1+1)
  +32(3\cdot3\cdot16+1)
  +10(32+1)
  =5{,}130.
\]
Writing this calculation before calling \texttt{count\_params()} turns the API
result into a checked prediction rather than an unexplained number.

\subsection{Pooling and receptive fields}

Max pooling reports the largest value in each local window; average pooling
reports the mean. Pooling reduces spatial resolution and can produce limited
local invariance, but it also discards information. Strided convolution is an
alternative learned downsampling mechanism.

The \emph{receptive field} of a unit is the region of the original input that
can influence it. Two stride-one $3\times3$ convolutions yield a $5\times5$
effective field; three yield $7\times7$. In general, if layer $\ell$ has
kernel $k_\ell$ and stride $s_\ell$, define the jump between adjacent units in
input coordinates by
\[
 j_0=1,
 \qquad j_\ell=j_{\ell-1}s_\ell,
\]
and receptive-field width by
\[
 r_0=1,
 \qquad r_\ell=r_{\ell-1}+(k_\ell-1)j_{\ell-1}.
\]
This recurrence exposes why pooling and stride make the field grow faster.

\begin{figure}[htbp]
\centering
\begin{tikzpicture}[
  stage/.style={draw=RiceBlue, rounded corners=1.5mm, minimum height=15mm,
    minimum width=19mm, align=center, font=\small},
  arrow/.style={-{Latex[length=2mm]}, thick, RiceBlue},
  node distance=4mm
]
\node[stage, fill=RiceBlue!10] (input) {image\\$28\times28\times1$};
\node[stage, fill=RiceGreen!15, right=of input] (conv1)
  {$3\times3$ conv\\16 maps};
\node[stage, fill=RiceGold!20, right=of conv1] (pool1)
  {$2\times2$ pool\\downsample};
\node[stage, fill=RiceGreen!15, right=of pool1] (conv2)
  {$3\times3$ conv\\32 maps};
\node[stage, fill=RiceGold!20, right=of conv2] (pool2)
  {$2\times2$ pool\\downsample};
\node[stage, fill=RiceBlue!10, right=of pool2] (head)
  {dense head\\10 logits};
\foreach \a/\b in {input/conv1,conv1/pool1,pool1/conv2,conv2/pool2,pool2/head}
  \draw[arrow] (\a) -- (\b);
\node[below=7mm of conv1, align=center, font=\scriptsize]
  {local shared detectors};
\node[below=7mm of conv2, align=center, font=\scriptsize]
  {larger effective field};
\node[below=7mm of head, align=center, font=\scriptsize]
  {task decision};
\end{tikzpicture}
\caption{A small convolutional classifier gradually trades spatial resolution
for more learned channels and a wider effective receptive field. Exact shapes
must be computed rather than inferred from the drawing.}
\label{fig:cnn-shape-story}
\end{figure}

\conceptcheck{A $28\times28$ input receives a $5\times5$ valid convolution
with stride two. Compute the output side length. Then compute the parameter
count for 16 filters when the input has one channel and when it has three.}

\section{From MNIST to cats and dogs}

MNIST is valuable because the data and computation are bounded. It is not a
template for every image claim. A real cats-versus-dogs study introduces
nuisance variables: backgrounds, photographers, compression, duplicate
animals, breed distribution, pose, occlusion, and label ambiguity. A model may
appear accurate after learning grass versus furniture.

A defensible pathway is:

\begin{enumerate}
  \item record source, license, checksum, and stable sample identifier;
  \item identify the independent sampling unit---preferably animal rather than
  image---and split groups before augmentation;
  \item validate decoding, shape, channel order, label vocabulary, and
  duplicates;
  \item establish a small CNN baseline under a bounded budget;
  \item compare transfer learning under the same split and selection rule;
  \item inspect confusion counts, per-class metrics, confidence, subgroup
  slices, and representative errors; and
  \item save and reload the selected artifact before reporting final test
  evidence.
\end{enumerate}

The official TensorFlow transfer-learning tutorial provides a reproducible
starting point for the mechanics.\footnote{TensorFlow, ``Transfer learning and
fine-tuning,''
\url{https://www.tensorflow.org/tutorials/images/transfer_learning}.} The
course notebook uses a tiny programmatic image fixture only to explain kernels.
It never presents the fixture as evidence about animal classification and does
not make CI depend on a large network download.

\begin{professionalbox}
Augmentation belongs inside the training boundary. Apply only transformations
that preserve the label in the intended operating domain. Never augment first
and then split: near-duplicate transformed images can leak across partitions.
Record the transform distribution because it is part of the model specification.
\end{professionalbox}

\section{Observable training}

Training is a running software process. When it behaves unexpectedly, a final
accuracy number cannot explain whether the failure came from data, numerical
optimization, the device pipeline, or the evaluation contract. Observability
provides structured evidence.

Python logging and TensorBoard serve complementary purposes:

\begin{description}
  \item[Application logs] record discrete events and context: configuration
  identity, data counts, warnings, exceptions, checkpoint paths, and final
  artifact identity.
  \item[TensorBoard summaries] record measurements indexed by step or time:
  losses, task metrics, learning rate, distributions, images, and selected
  performance traces.
\end{description}

Useful measurements have a diagnostic question attached:

\begin{center}
\begin{tabular}{p{.20\textwidth}p{.68\textwidth}}
\toprule
Boundary & Evidence to record \\
\midrule
Optimization & training loss, learning rate, gradient norm, nonfinite count \\
Generalization & validation loss, task metrics, train--validation gap \\
Data & class counts, corrupt records, observed ranges, augmentation rate \\
System & batch latency, examples per second, device/memory use when profiled \\
Release & frozen-test metrics, artifact hash, reload result, inference latency \\
\bottomrule
\end{tabular}
\end{center}

Metrics create overhead and risk. Histograms are larger than scalars;
profiling can materially slow execution; images can reveal private data; and
high-cardinality identifiers make logs expensive and difficult to aggregate.
Collect the least evidence that answers a defined diagnostic question.

\begin{workedexample}{diagnosing overfitting rather than naming it}
Training accuracy rises from $0.91$ to $0.999$, while validation loss reaches a
minimum in epoch four and then rises. TensorBoard makes the divergence visible,
but it does not prescribe a remedy. The team first verifies that the split unit
and preprocessing boundary are correct. It restores the epoch-four checkpoint,
compares class-specific errors, and repeats the experiment across seeds. Only
then does it evaluate regularization, augmentation, model capacity, or more
representative data. The plot begins an investigation; it does not complete one.
\end{workedexample}

The Keras TensorBoard callback writes event data at selected lifecycle points.
The normal full-MNIST route attaches this callback and records epoch-level loss
and accuracy for training and validation. The CI route deliberately attaches
no TensorBoard callback: it checks that the dependency and log-directory
configuration exist, but avoids creating an event writer whose startup and
filesystem costs could dominate a one-epoch smoke test. Different execution
budgets are legitimate only when their purposes and evidence claims remain
explicit.
The local viewer is launched from the repository root with
\begin{center}
\texttt{uv run tensorboard --logdir output/tensorboard/lecture-12 --port 6006}.
\end{center}
The viewer should bind locally unless exposure is explicitly authorized. Stop
it with \texttt{Ctrl+C}. A different port resolves a local collision. The
official Keras guide documents the callback and managed training lifecycle.\footnote{
TensorFlow, ``Training \& evaluation with the built-in methods,''
\url{https://www.tensorflow.org/guide/keras/training_with_built_in_methods}.}

\section{Debugging framework experiments}

Several failures recur often enough to treat them as tests:

\begin{enumerate}
  \item \textbf{Wrong kernel.} Installation succeeded in one interpreter but
  the notebook kernel points elsewhere. Inspect \texttt{sys.executable} and
  select Rice DSM.
  \item \textbf{Wrong rank or channel order.} Assert batch shapes at the data
  boundary and after the first batch.
  \item \textbf{Loss/output mismatch.} Test that a logits loss receives logits,
  and that sparse versus one-hot target encoding matches the selected loss.
  \item \textbf{Training/evaluation mode mismatch.} Dropout and batch
  normalization behave differently by mode. Reload in a fresh process and
  test inference explicitly.
  \item \textbf{Leakage.} Test the group split and ensure normalization or
  augmentation state is fitted only within the training boundary.
  \item \textbf{Empty telemetry.} Verify that the callback log directory and
  TensorBoard command resolve to the same absolute location.
  \item \textbf{Nondeterministic evidence.} Record seeds and deterministic
  options, then report variation across repeats rather than promising bitwise
  identity across all hardware.
\end{enumerate}

\begin{misconceptionbox}
A high validation accuracy is not an observability strategy. Validation
accuracy summarizes one outcome on one partition. Logs and telemetry help
explain the process that produced it; class-specific and operational evidence
determine whether it is fit for a decision.
\end{misconceptionbox}

\section{Exercises}
\begin{enumerate}
  \item Draw the framework stack in Figure~\ref{fig:tensorflow-stack} from
  memory and give one failure that belongs at each boundary.
  \item Derive the gradient of softmax cross-entropy with respect to each
  logit. Explain why a fused implementation is numerically preferable.
  \item Calculate every intermediate shape and parameter count in
  Figure~\ref{fig:cnn-shape-story} under valid and same padding.
  \item Prove the receptive-field recurrence and apply it to two
  convolution--pool blocks.
  \item Construct a $3\times3$ vertical-edge kernel and perform one complete
  patch calculation by hand. Explain why the operation is cross-correlation.
  \item Design a test that catches sparse versus one-hot loss mismatch before
  training begins.
  \item Compare a dense MNIST baseline with a small CNN under the same budget.
  Report parameter count, seed variation, error pairs, and runtime as well as
  accuracy.
  \item Design a group-aware cats-versus-dogs manifest. State which fields
  protect against duplication and leakage.
  \item Select no more than eight TensorBoard measurements for an MNIST run.
  For each, state the failure hypothesis it helps test.
  \item Extend the experiment to Fashion-MNIST. Frame a question about visual
  ambiguity between garment classes, then write a model-card limitation based
  on observed errors.
\end{enumerate}

\begin{chapterreview}
TensorFlow supplies tensor execution, mutable variables, automatic
differentiation, and optimized device kernels. Keras organizes these
capabilities into layers, models, objectives, optimizers, metrics, callbacks,
and a managed lifecycle. Concise framework code does not make the modeling
argument concise: shapes, label semantics, split boundaries, checkpoints, and
artifacts remain explicit contracts. Convolution fits image data by combining
local connectivity with parameter sharing. Channels represent learned feature
maps, stride and padding determine spatial shape, pooling downsamples, and
composition expands receptive fields. Python logging and TensorBoard make a
running experiment diagnosable when each measurement is tied to a concrete
question. They supplement rather than replace held-out decision evidence.
\end{chapterreview}

\begin{notebooklink}
Lecture 12 notebook 03 introduces the TensorFlow/Keras API, verifies automatic
differentiation against a hand calculation, and downloads the full official
MNIST dataset through Keras for the normal student analysis. CI explicitly
selects the repository's bounded authentic-MNIST sample as a software smoke
test. The notebook visualizes convolution mechanics, gives exact local
TensorBoard commands, and develops a real-data cats-versus-dogs extension with
provenance, licensing, and group-aware splitting requirements.
\end{notebooklink}
